\section{The Completed Determinant and Cyclotomic Twists}
\label{sec:completion-twists}

\subsection*{Base field and smooth projective models}

Throughout the geometric discussion we take \(k=\QQ\). Since
irreducibility, smoothness, and genus are unchanged after extending the
constant field to \(\overline{\QQ}\), we pass freely between \(k\) and
\(\overline{\QQ}\) when convenient. For an affine curve
\(C_{\mathrm{aff}}\subset\mathbb{A}^2_k\) defined by an irreducible
polynomial equation, we write \(\overline{X}\) (resp.\ \(\overline{Y}\))
for the smooth projective curve whose function field is
\(k(C_{\mathrm{aff}})\) (resp.\ its degree-two subfield under the
relevant involution). All genus and Hurwitz statements are made for
these complete nonsingular models.

\subsection{\texorpdfstring{The \(L1\) Zeckendorf normalisation interface}
  {The L1 Zeckendorf normalisation interface}}
\label{subsec:l1-zeckendorf-normalizer}

\ifshortcicm
This subsection records the integer normalisation protocol used in the
background comparison construction; no determinant proof below depends
on this protocol.

Let \(F_0=1\), \(F_1=2\), and \(F_{i+2}=F_{i+1}+F_i\).  A raw word is a
finite integer vector \(a=(a_0,\ldots,a_n)\), printed as
\([a_n,\ldots,a_0]\), with value
\[
  \operatorname{Val}(a)=\sum_{i=0}^n a_iF_i .
\]
A normal word is \(0\), \(+b_N\cdots b_0\), or \(-b_N\cdots b_0\),
where \(b_i\in\{0,1\}\), \(b_N=1\), and no two adjacent \(b_i\)'s are
both \(1\).  The operation \(\operatorname{NF}_\infty\) maps a raw word
to the signed Zeckendorf expansion of its value, using the greedy
algorithm.  Addition, multiplication, and the two digit-shift maps
\(\Phi_0,\Phi_1\) are defined by exact integer evaluation followed by
\(\operatorname{NF}_\infty\).

\begin{proposition}[Totality and size bounds for the \(L1\) normaliser]
\label{prop:l1-zeckendorf-normalizer-total}
The operations
\(\operatorname{NF}_\infty,\operatorname{Val},\operatorname{Add},
\operatorname{Mul},\Phi_0,\Phi_1\) are total on the grammars above.
If a raw input has \(n+1\) coefficients of absolute value at most
\(A\ge1\), its normal output has length \(O(n+\log A+1)\).  For normal
inputs of lengths \(p,q\), addition and \(\Phi_\epsilon\) have output
length \(O(\max\{p,q\}+1)\), while multiplication has output length
\(O(p+q)\).
\end{proposition}

\begin{proof}
Zeckendorf's theorem gives the unique non-adjacent binary expansion for
each non-negative integer \cite{Zeckendorf1972}; the external sign gives
the signed form.  Since \(F_i\asymp\varphi^i\), a raw value bounded by
\((n+1)AF_n\) needs only \(O(n+\log A+1)\) places.  The same exponential
growth estimate gives the stated bounds for addition, digit shifts, and
multiplication.  All operations use finite exact integer arithmetic, so
they are total.
\end{proof}

\else
The finite-state kernel below is the only object used in the determinant
proofs.  When the surrounding \(L1\) protocol is mentioned, however, we
use the following explicit Zeckendorf normalisation interface.  It is a
software-grade exact-integer protocol: every input and output has a
grammar, every operation is total, and the size bounds are polynomial in
the printed input.  It is not an assertion that the implementation uses
the best online or finite-automaton Zeckendorf arithmetic.  Such
efficient arithmetic is a separate subject; see
\cite{AhlbachUsatineFrougnyPippenger2013} for linear-time addition and
related circuit bounds.

Let \(\mathcal F_1=1\), \(\mathcal F_2=2\), and
\(\mathcal F_{s+2}=\mathcal F_{s+1}+\mathcal F_s\).  Thus
\(\mathcal F_1,\mathcal F_2,\mathcal F_3,\ldots=1,2,3,\ldots\).
For the printed digit strings it is convenient to use zero-based
places, so throughout this subsection
\[
  F_i:=\mathcal F_{i+1}\qquad(i\ge0);
\]
equivalently \(F_0=1\), \(F_1=2\), and
\(F_{i+2}=F_{i+1}+F_i\).  A raw \(L1\) word is a finite integer vector
\[
  a=(a_0,\ldots,a_n)\in\ZZ^{n+1},
\]
printed as a bracketed decimal list
\([a_n,\ldots,a_0]\); trailing zero coefficients may be suppressed, but
the empty list is not used.  Its value is
\[
  \operatorname{Val}(a):=\sum_{i=0}^n a_iF_i\in\ZZ .
\]
A normal \(L1\) word is one of the following strings:
\[
  0,\qquad +\,b_N\cdots b_0,\qquad -\,b_N\cdots b_0,
\]
where \(N\ge0\), \(b_i\in\{0,1\}\), \(b_N=1\), and
\(b_ib_{i+1}=0\) for every \(i\).  Its value is \(0\) in the first
case and
\[
  \pm\sum_{i=0}^N b_iF_i
\]
in the signed cases.  This is the usual Zeckendorf grammar
\cite{Zeckendorf1972}, with an explicit external sign.

\begin{lemma}[Fibonacci reindexing for imported interval maxima]
\label{lem:fibonacci-reindexing}
Let \((E_t)_{t\ge0}\) denote the Fibonacci convention
\[
  E_0=E_1=1,\qquad E_{t+2}=E_{t+1}+E_t,
\]
and let \((C_t)_{t\ge0}\) denote the conventional convention
\[
  C_0=0,\qquad C_1=1,\qquad C_{t+2}=C_{t+1}+C_t.
\]
Then, for every \(s\ge1\),
\[
  \mathcal F_s=E_s=C_{s+1}.
\]
Consequently an imported interval-maximum statement on
\([E_s,E_{s+1})\), or equivalently on \([C_{s+1},C_{s+2})\), is the
same statement on
\([\mathcal F_s,\mathcal F_{s+1})\).  In zero-based digit positions this
interval is
\[
  [\mathcal F_s,\mathcal F_{s+1})=[F_{s-1},F_s).
\]
\end{lemma}

\begin{proof}
The boundary cases are the only possible source of ambiguity.  For
\(s=1\),
\[
  \mathcal F_1=1=E_1=C_2,
\]
while for \(s=2\),
\[
  \mathcal F_2=2=E_2=C_3.
\]
For \(s\ge1\), the sequences
\((\mathcal F_s)_{s\ge1}\), \((E_s)_{s\ge1}\), and
\((C_{s+1})_{s\ge1}\) have the same first two terms and satisfy the same
recursion, so induction gives
\(\mathcal F_s=E_s=C_{s+1}\) for all \(s\ge1\).  Replacing the two
endpoints in any half-open interval by these equal quantities gives the
interval conversion.  Finally \(F_i=\mathcal F_{i+1}\), so
\(\mathcal F_s=F_{s-1}\) and \(\mathcal F_{s+1}=F_s\), giving the
zero-based digit-position form.
\end{proof}

\begin{center}
\fbox{\begin{minipage}{0.94\textwidth}
\textbf{\(L1\) normalisation algorithms.}

\begin{enumerate}[leftmargin=2.2em]
  \item \(\operatorname{Val}(x)\).  If \(x\) is raw, return
  \(\sum_i a_iF_i\).  If \(x\) is normal, return the signed
  Zeckendorf sum above.

  \item \(\operatorname{NF}_\infty(x)\).  Compute
  \(N=\operatorname{Val}(x)\).  If \(N=0\), return \(0\).  Otherwise
  record the sign of \(N\), replace \(N\) by \(|N|\), compute
  \(F_0,\ldots,F_K\) up to the largest \(F_K\le N\), and run the greedy
  loop:
  choose the largest unused \(F_j\le N\), set \(b_j=1\), replace
  \(N\leftarrow N-F_j\), and continue with indices at most \(j-2\).
  Return the recorded sign together with \(b_K\cdots b_0\), deleting
  leading zeroes.

  \item \(\operatorname{Add}(x,y)\).  Return
  \(\operatorname{NF}_\infty(\operatorname{Val}(x)+\operatorname{Val}(y))\).
  Equivalently, after padding raw coefficient vectors, add them
  coefficientwise and normalise.

  \item \(\operatorname{Mul}(x,y)\).  Return
  \(\operatorname{NF}_\infty(\operatorname{Val}(x)\operatorname{Val}(y))\).

  \item \(\Phi_0(x)\) and \(\Phi_1(x)\).  Let
  \(\operatorname{digits}(\operatorname{NF}_\infty(x))=b_N\cdots b_0\)
  be the unsigned magnitude digits, with all \(b_i=0\) for \(x=0\).
  For \(\epsilon\in\{0,1\}\), form the raw shifted word
  \[
    c_0=\epsilon,\qquad c_{i+1}=b_i\quad(0\le i\le N),
  \]
  and return \(\operatorname{NF}_\infty(c)\), with the original sign
  reapplied when \(x\ne0\).  Thus \(\Phi_\epsilon\) is the total
  normalised operation of appending the Fibonacci digit \(\epsilon\) on
  the least-significant side after shifting the old magnitude one
  Fibonacci place.
\end{enumerate}
\end{minipage}}
\end{center}

\begin{proposition}[Executable Zeckendorf protocol]
\label{prop:executable-zeckendorf-protocol}
Let a finite binary word \(c=(c_0,\ldots,c_N)\) be admissible when
\(c_k\in\{0,1\}\) and \(c_kc_{k+1}=0\) for all \(k\).  Put
\[
  \operatorname{Val}(c):=\sum_{k=0}^N c_kF_k .
\]
For \(n\in\NN\), define \(\operatorname{NF}_\infty(n)\), and for
admissible words \(c,d\), define
\(\operatorname{Add}(c,d)\), \(\operatorname{Mul}(c,d)\),
\(\Phi_0\), and \(\Phi_1\), by the following executable protocol.
\[
\begin{array}{ll}
\textsc{NF}_\infty(n):&
  r\leftarrow n,\quad c_k\leftarrow0\text{ for all computed }k;\\
& \textbf{while } r>0 \textbf{ do}\\
& \quad \text{choose the largest }F_k\le r;\\
& \quad c_k\leftarrow1,\quad r\leftarrow r-F_k;\\
& \textbf{return }c;\\[2mm]
\operatorname{Add}(c,d):&
  \textbf{return }\operatorname{NF}_\infty
  \bigl(\operatorname{Val}(c)+\operatorname{Val}(d)\bigr);\\[1mm]
\operatorname{Mul}(c,d):&
  \textbf{return }\operatorname{NF}_\infty
  \bigl(\operatorname{Val}(c)\operatorname{Val}(d)\bigr);\\[1mm]
\Phi_0(c_n)(f):& \textbf{return } f^{\circ n};\\
\Phi_1(c_n)(T):& \textbf{return } T^{\circ n},
\end{array}
\]
where \(c_n=\operatorname{NF}_\infty(n)\), \(f\) is an endomorphism,
\(T\) is an endomorphism operator, and \(h^{\circ0}\) denotes the
identity.  Then
\[
  \operatorname{Val}(\operatorname{NF}_\infty(n))=n,
\]
\[
  \operatorname{Val}(\operatorname{Add}(c,d))
    =\operatorname{Val}(c)+\operatorname{Val}(d),\qquad
  \operatorname{Val}(\operatorname{Mul}(c,d))
    =\operatorname{Val}(c)\operatorname{Val}(d),
\]
and, for \(c_a=\operatorname{NF}_\infty(a)\) and
\(c_b=\operatorname{NF}_\infty(b)\),
\[
  \Phi_0(c_{a+b})(f)=\Phi_0(c_a)(f)\circ\Phi_0(c_b)(f),\qquad
  \Phi_1(c_{ab})=\Phi_1(c_a)\circ\Phi_1(c_b).
\]
The first composition is taken in the endomorphism monoid containing
\(f\); the second is composition of the maps \(T\mapsto T^{\circ n}\)
on endomorphism operators.  Equivalently, after fixing \(f\), the first
identity reads
\[
  \Phi_0(c_{a+b})=\Phi_0(c_a)\circ\Phi_0(c_b)
\]
inside that endomorphism monoid.
This is a computable protocol, not an efficient online Zeckendorf
arithmetic algorithm.
\end{proposition}

\begin{proof}
The greedy loop terminates because each iteration subtracts a positive
Fibonacci weight from the remaining integer \(r\).  The usual
Zeckendorf uniqueness theorem \cite{Zeckendorf1972} says that each
non-negative integer has a unique admissible binary expansion in the
weights \(F_k\).  The greedy construction is precisely that expansion;
hence
\[
  \operatorname{Val}(\operatorname{NF}_\infty(n))=n.
\]
Applying this value identity to the definitions of addition and
multiplication gives
\[
  \operatorname{Val}(\operatorname{Add}(c,d))
    =\operatorname{Val}\!\left(
      \operatorname{NF}_\infty(\operatorname{Val}(c)+\operatorname{Val}(d))
      \right)
    =\operatorname{Val}(c)+\operatorname{Val}(d),
\]
and the same calculation gives
\[
  \operatorname{Val}(\operatorname{Mul}(c,d))
    =\operatorname{Val}(c)\operatorname{Val}(d).
\]
Finally, \(c_a\), \(c_b\), \(c_{a+b}\), and \(c_{ab}\) are only normal
names for the integers \(a\), \(b\), \(a+b\), and \(ab\).  Therefore,
for any fixed endomorphism \(f\),
\[
  \Phi_0(c_{a+b})(f)=f^{\circ(a+b)}
    =f^{\circ a}\circ f^{\circ b}
    =\Phi_0(c_a)(f)\circ\Phi_0(c_b)(f).
\]
Likewise, for any endomorphism operator \(T\),
\[
  ((\Phi_1(c_a)\circ\Phi_1(c_b))(T))
    =\Phi_1(c_a)(T^{\circ b})
    =\bigl(T^{\circ b}\bigr)^{\circ a}
    =T^{\circ ab}
    =\Phi_1(c_{ab})(T),
\]
with the convention that the zeroth iterate is the identity.  The
protocol uses exact integer evaluation followed by greedy normalisation;
it is therefore computable.  The claim is not that it achieves the
specialised online or linear-time bounds available for Zeckendorf
arithmetic, for which see
\cite{AhlbachUsatineFrougnyPippenger2013}.
\end{proof}

\begin{proposition}[Totality and size bounds for the \(L1\) normaliser]
\label{prop:l1-zeckendorf-normalizer-total}
The operations
\[
  \operatorname{NF}_\infty,\quad \operatorname{Val},\quad
  \operatorname{Add},\quad \operatorname{Mul},\quad \Phi_0,\quad \Phi_1
\]
defined above are total on their stated input grammars.  If the printed raw inputs have at most \(n+1\)
coefficients and each coefficient has absolute value at most \(A\ge1\),
then every normal output has length
\[
  O(n+\log A+1).
\]
For two normal inputs of lengths \(p\) and \(q\), the outputs of
\(\operatorname{Add}\), \(\Phi_0\), and \(\Phi_1\) have length
\(O(\max\{p,q\}+1)\), and the output of \(\operatorname{Mul}\) has
length \(O(p+q)\).  The algorithms run in time polynomial in the bit
length of the printed inputs using ordinary exact integer arithmetic.
\end{proposition}

\begin{proof}
The value map is a finite sum of integers, hence is total.  The greedy
part of \(\operatorname{NF}_\infty\) terminates because each selected
Fibonacci weight is positive and strictly decreases the remaining
non-negative integer.  It never selects adjacent weights: after choosing
\(F_j\), the algorithm restricts subsequent choices to indices at most
\(j-2\).  The standard Zeckendorf theorem gives existence and uniqueness
of the resulting non-adjacent binary expansion for every non-negative
integer \cite{Zeckendorf1972}; the external sign then gives a unique
normal form for every integer.  Thus \(\operatorname{NF}_\infty\) is
total and value preserving.  Since
\(\operatorname{Add}\), \(\operatorname{Mul}\), \(\Phi_0\), and
\(\Phi_1\) are finite compositions of \(\operatorname{Val}\), integer
arithmetic, finite coefficient shifts, and \(\operatorname{NF}_\infty\),
they are total as well.

It remains to record the size estimates.  The Fibonacci numbers satisfy
\(F_i\asymp\varphi^i\), where
\(\varphi=(1+\sqrt5)/2\).  For a raw vector with \(n+1\) coefficients
bounded by \(A\),
\[
  |\operatorname{Val}(a)|
  \le (n+1)A F_n
  \le C(n+1)A\varphi^n
\]
for an absolute constant \(C\).  Therefore the largest Fibonacci index
needed in the normal form is \(O(n+\log A+1)\).  For normal inputs of
lengths \(p\) and \(q\), their absolute values are
\(O(\varphi^p)\) and \(O(\varphi^q)\).  Addition and the two
\(\Phi_\epsilon\) operations therefore need only
\(O(\max\{p,q\}+1)\) Fibonacci positions, while multiplication has
absolute value \(O(\varphi^{p+q})\) and hence needs \(O(p+q)\)
positions.  The displayed algorithms compute finitely many Fibonacci
weights and perform exact integer additions, multiplications,
comparisons, and divisions by subtraction/search over these lists; with
standard binary integer routines this is polynomial in the printed input
length.  These are auditability bounds, not the sharper streaming bounds
available for specialised Zeckendorf adders.
\end{proof}
\fi

\subsection{The unweighted spectrum}

\begin{proposition}[Explicit determinant]\label{prop:kernel-delta-explicit}
One has
\[
  \Delta(z,u)
  =1-(1+u)z-5uz^2+3u(1+u)z^3-u(u^2-3u+1)z^4
\]
\[
  \qquad
  +u(u^3-3u^2-3u+1)z^5+u^2(u^2+u+1)z^6.
\]
\end{proposition}

\begin{proof}
A symbolic determinant computation from the explicit \(10\times 10\)
matrix \(B(u)\) of \Cref{def:sync-kernel} gives the displayed formula;
the displayed polynomial is the exact certificate used in the sequel,
and the command block producing it is recorded in
\Cref{sec:appendix-computation}\,\textup{(a)}.
As a consistency check, the closed form satisfies the self-duality
identity of \Cref{prop:sync-self-duality}; in particular, comparing
\(\Delta(z,u)\) with \(\Delta(uz,u^{-1})\) pairs the coefficients of
\(z^k u^\ell\) and \(z^k u^{k-\ell}\) exactly as required.
\end{proof}

Although \(B(u)\) is a \(10\times 10\) matrix, the determinant has
degree \(6\) in \(z\) because four rows are redundant for every \(u\):
rows \(5,6,8,10\) are \(u\)-multiples of rows \(1,2,7,9\),
respectively. Hence \(\operatorname{rank} B(u)\le 6\), so \(0\) is a
persistent eigenvalue of multiplicity at least four.

\begin{proposition}[Untwisted Perron root and residue constant]
  \label{prop:sync-kernel-spectrum}
  The matrix \(B=B(1)\) is primitive and every row of \(B\) has sum
  \(3\). Consequently the spectral radius is determined before any
  characteristic-polynomial factorisation is used:
  \[
    B\mathbf 1=3\mathbf 1,\qquad
    \|B\|_\infty=\max_i\sum_j |B_{ij}|=3,\qquad
    \rho(B)\le \|B\|_\infty,
  \]
  so \(\rho(B)=3\), and the Perron eigenvalue \(3\) is simple and is
  the unique eigenvalue of modulus \(3\).
  In addition,
  \[
    \det(I-zB)=(z-1)(z+1)(3z-1)(z^3-z^2+z+1).
  \]
  The Perron residue constant \(C_B\) of
  \Cref{def:perron-residue-constant} is
  \[
    C_B=\frac{243}{272}.
  \]
\end{proposition}

\begin{proof}
Every state reaches \(\st{000}\) in at most three steps by the paths
listed in \Cref{def:sync-kernel}. Conversely, the transition table gives
explicit paths from \(\st{000}\) to the other states:
\[
  \st{000}\to\st{001},\qquad
  \st{000}\to\st{002},\qquad
  \st{000}\to\st{001}\to\st{010},
\]
\[
  \st{000}\to\st{001}\to\st{100},\qquad
  \st{000}\to\st{001}\to\st{101},\qquad
  \st{000}\to\st{002}\to\st{11-1},
\]
\[
  \st{000}\to\st{001}\to\st{010}\to\st{0-12},\qquad
  \st{000}\to\st{001}\to\st{010}\to\st{0-12}\to\st{01-1},
\]
\[
  \st{000}\to\st{002}\to\st{11-1}\to\st{1-12}.
\]
Thus the digraph of \(B\) is strongly connected. Since there is also the
self-loop \(\st{000}\to\st{000}\), the period is \(1\), so \(B\) is primitive.

We next determine the spectral radius directly from the non-negative
matrix \(B\), before using any determinant factorisation. Inspecting the
transition table at \(u=1\), every row of \(B=B(1)\) has sum \(3\).
Therefore
\[
  B\mathbf 1=3\mathbf 1,\qquad
  \|B\|_\infty=\max_i\sum_j |B_{ij}|=\max_i\sum_j B_{ij}=3,
  \qquad
  \rho(B)\le\|B\|_\infty.
\]
The first equality exhibits \(3\) as an eigenvalue. The inequality
follows from the \(\ell^\infty\)-operator norm, or equivalently from
\Cref{lem:row-sum-perron-bound}, and shows that no eigenvalue of \(B\)
can have modulus greater than \(3\). Hence \(3\le\rho(B)\le3\), so
\(\rho(B)=3\). Since \(B\) is primitive, the Perron--Frobenius theorem
gives that this eigenvalue is simple and that no other eigenvalue lies
on \(\abs{\lambda}=3\).

The determinant is used only after this Perron conclusion has been
proved, and only for the displayed factorisation and residue. Setting
\(u=1\) in
\Cref{prop:kernel-delta-explicit} gives
\[
  \det(I-zB)=(z-1)(z+1)(3z-1)(z^3-z^2+z+1).
\]
This factorisation is deliberately not used to prove the spectral
radius. In the variable \(z\), the cubic factor \(z^3-z^2+z+1\) has a
zero with \(\abs{z}<1\), so the false reciprocal-root argument would not
distinguish \(3z-1\) from all other factors. The Perron conclusion above
comes only from \(B\mathbf1=3\mathbf1\), the row-sum norm bound
\(\rho(B)\le\|B\|_\infty=3\), and primitivity. The factorisation is used
here solely for the residue computation
\[
  \begin{aligned}
  C_B
  &=\lim_{z\to1/3}
    \frac{1-3z}{(z-1)(z+1)(3z-1)(z^3-z^2+z+1)}  \\
  &=-\frac{1}{(-2/3)(4/3)(34/27)}
    =\frac{243}{272}.
  \end{aligned}
\]
This proves part~\textup{(i)} of
\Cref{thm:intro-kernel-geometry}.
\end{proof}

\ifshortcicm\else
Together with \Cref{lem:primitive-orbit-asymptotic}, this yields
\[
  p_n(B)=\frac{3^n}{n}
  +O\!\left(\frac{\max\{\Lambda(B)^n,3^{n/2}\}}{n}\right).
\]
\fi

\subsection{Completion and algebraic geometry}

Introducing
\[
  u=r^2,\qquad s=r+r^{-1},\qquad w=zr,
\]
the self-duality of \(\Delta(z,u)\) implies that
\(\Delta(w/r,r^2)\) is a polynomial in \(w\) and \(s\), which we denote
by \(\widehat\Delta(w,s)\).

\begin{proposition}[Explicit completed determinant]
  \label{prop:sync-hatdelta-completion}
  The completed determinant is
  \[
    \widehat\Delta(w,s)
    =1-sw-5w^2+3sw^3+(5-s^2)w^4+(s^3-6s)w^5+(s^2-1)w^6.
  \]
  It satisfies
  \[
    \widehat\Delta(w,-s)=\widehat\Delta(-w,s).
  \]
\end{proposition}

\begin{proof}
Substitute \(u=r^2\) and \(z=w/r\) into the explicit formula of
\Cref{prop:kernel-delta-explicit}. Terms in \(r\) and \(r^{-1}\)
combine to give a polynomial in \(s=r+r^{-1}\), and a direct expansion
produces the displayed formula. This displayed formula, together with
the parity identity below it, is the certificate used in the geometric
and cyclotomic arguments; the command block producing it is
recorded in \Cref{sec:appendix-computation}\,\textup{(b)}. The parity
relation is immediate. This proves part~\textup{(ii)} of
\Cref{thm:intro-kernel-geometry}.
\end{proof}

\begin{remark}
The phrase ``completed determinant'' is non-standard. In this note it
simply denotes the polynomial obtained by rewriting \(\Delta(w/r,r^2)\)
in the self-dual invariants \(w\) and \(s=r+r^{-1}\).
\end{remark}



\begin{proposition}[Generic Galois group]
  \label{prop:sync-hatdelta-galois-s6-generic}
  Viewed as a polynomial in \(w\), the completed determinant
  \(\widehat\Delta(w,s)\) has generic Galois group \(S_6\) over
  \(\QQ(s)\).
\end{proposition}

\begin{proof}
\ifshortcicm
Let \(G\le S_6\) be the generic Galois group over \(\QQ(s)\).  The
certificate runner computes the discriminant \(D(s)=\disc_w
\widehat\Delta(w,s)\) and prints
\((\deg D,\operatorname{lc}D)=(20,-256)\).  Hence \(D\) is not a square
in \(\QQ(s)\), so \(G\not\subset A_6\).

The same runner gives three separable specialisations and finite-field
factorisations.  At \(s=3\) modulo \(7\), the degree-six polynomial is
irreducible, giving a \(6\)-cycle in \(G\).  At \(s=2\) modulo \(7\),
the factorisation type is \((1)(1)(1)(1)(2)\), giving a transposition.
At \(s=3\) modulo \(19\), the factorisation type is \((1)(5)\), giving
a \(5\)-cycle.  Dedekind's theorem and the standard specialisation
theorem justify lifting these cycle types to the generic group
\cite{Serre1979,Serre2008Topics}.

A transitive subgroup of \(S_6\) containing a \(5\)-cycle is primitive:
any block system with \(2\) or \(3\) blocks would be fixed by the
5-cycle, forcing a block to contain its whole five-point support.
A primitive subgroup of \(S_6\) containing a transposition is \(S_6\)
\cite[Ch.~3, Theorem~3.3A]{Dixon1996}.  Thus \(G=S_6\).
\else
Let \(G\le S_6\) be the generic Galois group of
\(\widehat\Delta(w,s)\) over \(\QQ(s)\). We identify \(G\) in four
steps.

\medskip
\noindent\textit{Step 1 (discriminant, \(G\not\subset A_6\)).}\quad
Write \(D(s):=\disc_w\widehat\Delta(w,s)\in\QQ[s]\). The explicit
SageMath~10.2 computation recorded in
\Cref{sec:appendix-computation}\,(d) prints the output
\((20,-256)\) for \((\deg D,\operatorname{lc}D)\), so \(D(s)\) has
degree~\(20\) with leading coefficient~\(-256\).
This metadata is already a nonsquare certificate. Indeed, if a
polynomial \(D\in\QQ[s]\) is a square in \(\QQ(s)\), write
\(D=(A/B)^2\) with coprime \(A,B\in\QQ[s]\). Then
\(A^2=DB^2\), so every irreducible divisor of \(B\) divides \(A\);
hence \(B\) is constant and \(D\) is a square in \(\QQ[s]\). Its
leading coefficient is therefore a square in \(\QQ^\times\), in
particular positive. The leading coefficient \(-256\) rules this out.
Since the discriminant of \(\widehat\Delta(w,s)\) as a polynomial in
\(w\) changes sign under odd permutations,
\(D(s)\notin \QQ(s)^2\) implies \(G\not\subset A_6\).

\medskip
\noindent\textit{Step 2 (specialisation input).}\quad
We use the following standard fact: if \(s_0\in\QQ\) is such that
\(\widehat\Delta(w,s_0)\) is separable over \(\QQ\) (equivalently,
\(D(s_0)\neq 0\)), then the Galois group \(G_{s_0}\) of
\(\widehat\Delta(w,s_0)\) over \(\QQ\) embeds into \(G\); see the
specialisation theorem in \cite[\S 4.4, Proposition~1]{Serre2008Topics}.
The condition \(D(s_0)\neq 0\) is
verified at each chosen \(s_0\) by computing
\(\gcd(\widehat\Delta(w,s_0),\partial_w\widehat\Delta(w,s_0))\bmod p\)
for a convenient prime \(p\).

To read off cycle types inside \(G_{s_0}\), we apply Dedekind's
theorem in its localized form. Namely, let
\(f(w)\in\ZZ[w]\) have degree \(n\), and let \(p\) be a prime
which divides neither the leading coefficient of \(f\) nor
\(\disc(f)\). Then \(f\) has degree \(n\) after reduction modulo
\(p\), and multiplying its reduction by the inverse of the reduced
leading coefficient gives a monic polynomial with the same irreducible
factor degrees. If this reduction factors as
\(f_1\cdots f_r\) with the \(f_i\) distinct and irreducible over
\(\mathbb{F}_p\), then the Frobenius at \(p\) in \(\Gal(f/\QQ)\)
has cycle type \((\deg f_1,\dots,\deg f_r)\). This is the usual
Dedekind theorem \cite[I, \S{}6, Proposition~25]{Serre1979}, applied
over \(\ZZ_{(p)}\), in the Frobenius form used in
\cite[\S{}4.3]{Serre2008Topics}. The exact mod-\(p\) factorisation
data are recorded in \Cref{sec:appendix-computation}\,(d).

\medskip
\noindent\textit{Step 3 (three Frobenius specialisations).}

\smallskip
\noindent\textit{Transitivity (\(s=3\), mod \(7\)).}\quad
One checks that
\(\gcd(\widehat\Delta(w,3),\partial_w\widehat\Delta(w,3))\equiv 1\pmod 7\),
so \(D(3)\neq 0\) and the specialisation is separable.
Here \(\operatorname{lc}_w\widehat\Delta(w,3)=8\not\equiv0\pmod7\),
so the nonmonic normalization above is legitimate.
Modulo \(7\),
\[
  \widehat\Delta(w,3)
  \equiv w^6+2w^5+3w^4+2w^3+2w^2+4w+1\pmod 7,
\]
which is irreducible over \(\mathbb{F}_7\). Hence the Frobenius at \(7\)
in \(G_3\subset G\) is a \(6\)-cycle, so \(G\) acts transitively
on the six roots.

\smallskip
\noindent\textit{Transposition (\(s=2\), mod \(7\)).}\quad
One checks that
\(\gcd(\widehat\Delta(w,2),\partial_w\widehat\Delta(w,2))\equiv 1\pmod 7\),
so \(D(2)\neq 0\). Here
\(\operatorname{lc}_w\widehat\Delta(w,2)=3\not\equiv0\pmod7\).
Modulo \(7\),
\[
\begin{aligned}
  \widehat\Delta(w,2)
  &\equiv (w+6)(w+5)(w+2)(w+1) \\
  &\qquad\cdot(3w^2+3w+2)\pmod 7,
\end{aligned}
\]
The quadratic factor \(3w^2+3w+2\) is irreducible over
\(\mathbb{F}_7\). Thus the factorisation type is
\((1)(1)(1)(1)(2)\), so the Frobenius at \(7\) in \(G_2\subset G\) is a
transposition. Hence \(G\) contains a transposition.

\smallskip
\noindent\textit{5-cycle (\(s=3\), mod \(19\)).}\quad
One checks that
\(\gcd(\widehat\Delta(w,3),\partial_w\widehat\Delta(w,3))\equiv 1\pmod{19}\).
This squarefree reduction is exactly the separability witness:
equivalently \(19\nmid D(3)\), so Dedekind's theorem applies to this
factorisation.
Here \(\operatorname{lc}_w\widehat\Delta(w,3)=8\not\equiv0\pmod{19}\).
Modulo \(19\),
\[
  \widehat\Delta(w,3)
  \equiv (w+16)\bigl(8w^5+14w^4+9w^2+3w+6\bigr)\pmod{19},
\]
with the quintic factor irreducible over \(\mathbb{F}_{19}\). The
factorisation type is \((1)(5)\), so the Frobenius at \(19\) in
\(G_3\subset G\) is a \(5\)-cycle. Hence \(G\) contains a
\(5\)-cycle.

\medskip
\noindent\textit{Step 4 (subgroup classification).}\quad
We now have: \(G\le S_6\) is transitive, contains a transposition, and
contains a \(5\)-cycle. A transitive subgroup of \(S_6\) that contains
a \(5\)-cycle is primitive. Indeed, a non-trivial block system on six
points has either \(2\) or \(3\) blocks. Let \(\sigma\in G\) be a
\(5\)-cycle. The induced action of \(\sigma\) on such a block system
has order dividing both \(5\) and the order of a permutation group on
\(2\) or \(3\) points; hence the induced action is trivial. Therefore
each block is fixed setwise by \(\sigma\). Any block meeting the support
of the \(5\)-cycle must then contain the whole five-point support,
which is impossible for block size \(3\) or \(2\). Thus no non-trivial
block system exists, so \(G\) is primitive
\cite[Ch.~3, \S{}3.3]{Dixon1996}.
A primitive subgroup of \(S_6\) containing a transposition is the full
symmetric group \cite[Ch.~3, Theorem~3.3A]{Dixon1996}. Hence \(G=S_6\),
completing part~\textup{(iii)} of \Cref{thm:intro-kernel-geometry}.
\fi
\end{proof}

\begin{corollary}\label{lem:hatdelta-irreducible}
The polynomial \(\widehat\Delta(w,s)\) is irreducible in \(\QQ[w,s]\).
\end{corollary}

\begin{proof}
Since \(\Gal(\widehat\Delta/\QQ(s))=S_6\), the action on the six roots is
transitive, so \(\widehat\Delta(w,s)\) is irreducible over \(\QQ(s)\).
Gauss's lemma then implies irreducibility in \(\QQ[w,s]\).
\end{proof}
We first construct the quotient curve \(\overline{Y}\) and establish its
smoothness; then we record the integrality of \(C_{\mathrm{aff}}\) and
the non-squareness of \(x\) in \(k(\overline{Y})\), which together
justify the definition of \(\overline{X}\) as a well-defined smooth
projective curve.

\begin{proposition}[The quotient is a smooth plane quartic]
  \label{prop:sync-hatdelta-quotient}
  The invariants
  \[
    x:=w^2,\qquad y:=sw
  \]
  identify the sign quotient birationally with the projective closure of
  the affine curve
  \[
    F(x,y):=
    1-y-5x+3xy+5x^2-xy^2+xy^3-6x^2y+x^2y^2-x^3=0.
  \]
  This projective closure is smooth; consequently it is the smooth
  projective model \(\overline{Y}\), and
  \[
    g(\overline{Y})=3.
  \]
\end{proposition}

\begin{proof}
\ifshortcicm
\Cref{thm:sign-quotient-plane-genus} applies because
\Cref{lem:self-dual-compression} gives the required parity bounds and
\(\widehat\Delta(0,s)=1\).  Substituting \(x=w^2\), \(y=sw\) gives the
displayed quartic \(F\).  The affine smoothness certificate is the
Groebner basis
\(\langle F,\partial_xF,\partial_yF\rangle=\langle1\rangle\), printed
by the runner.  At infinity, the homogenised leading form is
\(XY^2(X+Y)\), giving \([0:1:0]\), \([1:0:0]\), and \([1:-1:0]\), with
nonzero partials \(1\), \(-1\), and \(4\).  Thus the projective quartic
is smooth and has genus \((4-1)(4-2)/2=3\).
\else
By \Cref{lem:self-dual-compression}, the coefficient of \(w^j\) in
\(\widehat\Delta(w,s)\) has degree at most \(j\) in \(s\) and the
parity of \(j\). Hence
\Cref{thm:sign-quotient-plane-genus} applies to the involution
\(\iota:(w,s)\mapsto(-w,-s)\): the invariants
\[
  x=w^2,\qquad y=sw
\]
identify the invariant function field birationally, and the theorem
gives a unique polynomial \(F(x,y)\in\QQ[x,y]\) with
\(\widehat\Delta(w,s)=F(w^2,sw)\).
A direct substitution gives the displayed formula for \(F\). Since
\(\widehat\Delta(0,s)=1\), the constant coefficient hypothesis
\(c_0\ne0\) in \Cref{thm:sign-quotient-plane-genus} is satisfied; hence
the quotient function field is the function field of \(F=0\). Thus
\(\overline{Y}\) is birational to the projective closure of the affine
quartic \(\{F=0\}\).
After the smoothness check below, this closure is the smooth projective
model \(\overline{Y}\).

We verify that this projective closure is smooth. Homogenizing
\(F(x,y)\) to degree \(4\) in projective coordinates \([X:Y:Z]\) gives
\[
\begin{aligned}
  F_h(X,Y,Z)
  &= Z^4 - YZ^3 - 5XZ^3 + 3XYZ^2 + 5X^2Z^2 - XY^2Z \\
  &\quad + XY^3 - 6X^2YZ + X^2Y^2 - X^3Z .
\end{aligned}
\]
On the affine chart \(Z=1\),
\[
  \partial_x F
  =-5+3y+10x-y^2+y^3-12xy+2xy^2-3x^2,
\]
\[
  \partial_y F
  =-1+3x-2xy+3xy^2-6x^2+2x^2y.
\]
A Gr\"obner-basis computation of the ideal
\(\langle F,\,\partial_x F,\,\partial_y F\rangle\subset\QQ[x,y]\)
gives \(\langle 1\rangle\), so there is no affine singular point.

The line at infinity \(Z=0\) meets the quartic where the leading form
\(XY^2(X+Y)=0\), giving the three points
\[
  [0:1:0],\qquad [1:0:0],\qquad [1:-1:0].
\]
At \([0:1:0]\) (chart \(Y=1\)): one computes
\[
\begin{aligned}
  \partial_X F_h\big|_{(0,1,0)}
  &=
  \bigl[-5Z^3+3YZ^2+10XZ^2-Y^2Z+Y^3 \\
  &\qquad -12XYZ+2XY^2-3X^2Z\bigr]_{(0,1,0)}
  = 1\neq 0 .
\end{aligned}
\]
At the other two points at infinity,
\[
  \partial_Z F_h\big|_{(1,0,0)}=-1\neq 0,
  \qquad
  \partial_Z F_h\big|_{(1,-1,0)}=4\neq 0 .
\]

Hence the projective curve is smooth everywhere, so by the
\Cref{thm:sign-quotient-plane-genus} with \(n=4\), this nonsingular
quartic is the smooth projective quotient model and
\[
  g(\overline{Y})=(4-1)(4-2)/2=3.
\]
\fi
\end{proof}

\begin{lemma}[Integrality and function-field setup]
  \label{lem:integral-function-field}
  Let \(k=\QQ\) and let \(\overline{Y}\) be the smooth projective quartic
  of \Cref{prop:sync-hatdelta-quotient}. Then:
  \begin{enumerate}[label=\textup{(\roman*)}]
    \item The affine curve \(C_{\mathrm{aff}}:=\{\widehat\Delta(w,s)=0\}
          \subset\mathbb{A}^2_k\) is integral.
    \item The element \(x=w^2\in k(\overline{Y})\) is not a square
          in \(k(\overline{Y})\).
    \item The extension \(k(\overline{Y})(\sqrt{x})/k(\overline{Y})\)
          has degree exactly two, and its function field is \(k(C_{\mathrm{aff}})\).
  \end{enumerate}
\end{lemma}

\begin{proof}
\ifshortcicm
Irreducibility follows from the generic \(S_6\) group, hence the affine
curve is integral.  On the quotient quartic, \(F(0,y)=1-y\), so
\(x=0\) meets at \(P=(0,1)\), and
\(\partial_yF(0,1)=-1\) gives \(\operatorname{ord}_P(x)=1\).  Thus
\(x\) is not a square in \(k(\overline{Y})\).  Adjoining
\(\sqrt{x}=w\), with \(s=y/w\), recovers the completed curve function
field.
\else
Part~(i): \(\widehat\Delta(w,s)\) is irreducible in \(\QQ[w,s]\) by
\Cref{lem:hatdelta-irreducible}, so its zero locus is an irreducible
variety. Since \(\widehat\Delta\) is prime in the UFD \(\QQ[w,s]\), the
coordinate ring \(\QQ[w,s]/(\widehat\Delta)\) is a domain, so the
variety is also reduced. Hence \(C_{\mathrm{aff}}\) is integral.

Part~(ii): From \Cref{prop:sync-hatdelta-quotient}, the affine quartic
satisfies \(F(0,y)=1-y\), so \(x=0\) meets \(\overline{Y}\) only at
the point \(P=(0,1)\), with \(\partial_y F(0,1)=-1\neq 0\), giving
\(\operatorname{ord}_P(x)=1\). A square in \(k(\overline{Y})\) would
have even valuation at every place; the odd order at \(P\) shows
\(x\notin k(\overline{Y})^2\).

Part~(iii): Since \(x\notin k(\overline{Y})^2\), the extension
\(k(\overline{Y})(\sqrt{x})/k(\overline{Y})\) has degree two. The
identifications \(w=\sqrt{x}\) and \(s=y/\sqrt{x}\) show that this
extension is the function field of \(C_{\mathrm{aff}}\).
\fi
\end{proof}

\begin{proposition}[Two-sheeted cover of genus six]
  \label{prop:sync-hatdelta-double-cover}
  The covering map \(\overline{X}\to\overline{Y}\) is a degree-two map
  branched at exactly two points of \(\overline{Y}\). Hence
  \[
    g(\overline{X})=6.
  \]
\end{proposition}

\begin{proof}
\ifshortcicm
The cover is the quadratic extension
\(k(\overline{Y})(\sqrt{x})/k(\overline{Y})\).  The divisor calculation
on the quartic gives
\[
  \operatorname{div}(x)=P+2Q_1-2Q_0-Q_\infty,
\]
so the odd valuations, hence the ramification points, are exactly
\(P\) and \(Q_\infty\).  Riemann--Hurwitz then gives
\[
  2g(\overline{X})-2=2(2\cdot3-2)+2,
\]
and \(g(\overline{X})=6\).
\else
By \Cref{lem:integral-function-field}, the extension of function fields
\[
  k(\overline{X})/k(\overline{Y})
\]
is obtained by adjoining \(\sqrt{x}\), and it has degree exactly two.

We use the standard criterion for ramification in quadratic extensions
of function fields: let \(K\) be a function field of characteristic
\(\neq 2\), and let \(L=K(\sqrt{f}\,)\) with \(f\in K^\times\) not a
square. Then a place \(P\) of \(K\) ramifies in \(L/K\) if and only if
\(\operatorname{ord}_P(f)\) is odd; see \cite[III.7.3]{Stichtenoth2009}.
Here \(f=x\), and we determine the divisor of \(x\) on \(\overline{Y}\).

On the affine chart, \(F(0,y)=1-y\), so the only affine point with
\(x=0\) is \(P=(0,1)\). Since \(\partial_y F(0,1)=-1\neq 0\), the line
\(x=0\) meets \(\overline{Y}\) transversely there, and therefore
\[
  \operatorname{ord}_P(x)=1.
\]

For the points at infinity, work on the chart \(X=1\) with
\(y=Y/X\) and \(z=Z/X\). The equation becomes
\[
  G(y,z)=z^4-yz^3-5z^3+3yz^2+5z^2-y^2z+y^3-6yz+y^2-z.
\]
At \(Q_0=[1:0:0]\), one has \(\partial_z G(0,0)=-1\), so \(z\) is an
analytic function of \(y\), and solving \(G(y,z)=0\) gives
\[
  z=y^2-5y^3+O(y^4).
\]
Hence \(x=X/Z=1/z\) has pole order \(2\) at \(Q_0\), so
\[
  \operatorname{ord}_{Q_0}(x)=-2.
\]
At
\(Q_\infty=[1:-1:0]\), writing \(y=-1+t\) gives
\[
  G(-1+t,z)=t-2t^2+t^3+4z-4tz-t^2z
  +2z^2+3tz^2-4z^3-tz^3+z^4,
\]
so \(\partial_z G(-1,0)=4\neq 0\), and
\[
  z=-\frac14t+\frac{7}{32}t^2+O(t^3).
\]
Thus \(x=1/z\) has a simple pole at \(Q_\infty\), so
\[
  \operatorname{ord}_{Q_\infty}(x)=-1.
\]

At the remaining point \(Q_1=[0:1:0]\), use the chart \(Y=1\) with
\(u=X/Y\) and \(v=Z/Y\). The homogenised equation begins as
\[
  u-v^3+O(uv,v^4,u^2)=0,
\]
so \(u=v^3+O(v^4)\), and therefore \(x=X/Z=u/v\) has a zero of order
\(2\) at \(Q_1\), so
\[
  \operatorname{ord}_{Q_1}(x)=2.
\]

Thus
\[
  \operatorname{div}(x)=P+2Q_1-2Q_0-Q_\infty.
\]
The odd valuations occur exactly at \(P\) and \(Q_\infty\). By the
quadratic ramification criterion, these are exactly the two ramified
places of \(k(\overline{X})/k(\overline{Y})\), and each contributes one
simple branch point. The Riemann--Hurwitz genus formula
\cite[III.4.12]{Stichtenoth2009} then gives
\[
  2g(\overline{X})-2=2(2g(\overline{Y})-2)+2=2(2\cdot 3-2)+2,
\]
hence \(g(\overline{X})=6\).
\fi
\end{proof}





\subsection{Cyclic lifts and primitive real factors}

For \(q\ge 1\), let \(P_q\) be the cyclic permutation matrix on
\(\ZZ/q\ZZ\), acting on the standard basis by \(P_qe_t=e_{t+1}\). The
natural \(q\)-cyclic skew-product lift attached to the binary output
labelling is
\[
  \widetilde B_q:=B_0\otimes I_q+B_1\otimes P_q.
\]
Its state set is \(\mathcal S\times\ZZ/q\ZZ\), and an edge
\(x\to y\) of output \(e\in\{0,1\}\) in the original kernel becomes
\((x,t)\to(y,t+e)\) in the lift. This is the standard cyclic extension
attached to the output cocycle
\cite{AdlerKitchensMarcus1985,Fiebig1993,BoyleSchmieding2017}.

\begin{definition}[Admissible cyclic covers]
  \label{def:admissible-cyclic-cover}
  Fix a finite non-negative binary-output kernel
  \(B(u)=B_0+uB_1\).  An admissible cyclic cover of level \(q\) is the
  \(\ZZ/q\ZZ\)-skew product obtained from the displayed output cocycle:
  its adjacency operator is
  \[
    B_0\otimes I_q+B_1\otimes P_q,
  \]
  so an output-\(e\) edge advances the fibre coordinate by \(e\) and no
  additional vertices, edges, weights, or cocycles are inserted.  The
  admissible cover category for \(B\) has these covers as objects and
  the evident fibre translations and level maps induced by quotients
  \(\ZZ/q'\ZZ\to\ZZ/q\ZZ\) as morphisms.

  For the ten-state kernel of \Cref{def:sync-kernel}, the admissible
  objects are exactly the lifts represented by the matrices
  \[
    \widetilde B_q=B_0\otimes I_q+B_1\otimes P_q.
  \]
  The printed matrix and certificate data determine this category.  The
  matrices \(B_0\) and \(B_1\) are read from the transition table and
  from the displayed matrix \(B(u)\) in
  \Cref{sec:appendix-computation}.  The determinant, self-duality,
  endpoint-separation, and finite-\(m\) branch-selection checks are
  named outputs of the Sage runner.  Concrete
  certificate examples are the primitive real factors for
  \(m=4,6,8,10\) displayed in
  \Cref{sec:appendix-computation}\,\textup{(e)} and the exact
  finite-\(m\) branch-selection check for \(4\le m<32\); each is
  obtained from an admissible cyclic lift at level \(m\).  Thus the
  covers used below are not arbitrary auxiliary covers; they are the
  canonical cyclic covers forced by the certified binary output
  labelling.
\end{definition}

\begin{remark}[Scope of the cover-relative statements]
  \label{rem:cover-scope}
  The obstruction and branch-selection statements below are relative to
  the admissible cyclic-cover category of
  \Cref{def:admissible-cyclic-cover}.  One may attach unrelated
  auxiliary covers to the same underlying graph, for example by adding
  a one-spike voltage feature or by changing the cocycle on a selected
  edge.  Such constructions can be useful stress tests for how
  determinant certificates behave under perturbation, but they are not
  intrinsic behaviour of the displayed self-dual kernel and they are
  not part of the software-certificate claim made in this paper.
\end{remark}

\begin{proposition}[Fourier decomposition of the cyclic lift]
  \label{prop:sync-cyclic-lift-factorisation}
  Let \(\omega_q=e^{2\pi i/q}\). Then
  \[
    \det(I-z\widetilde B_q)
    =\prod_{j=0}^{q-1}\Delta(z,\omega_q^j).
  \]
\end{proposition}

\begin{proof}
Let \(F_q=(\omega_q^{ab})_{0\le a,b\le q-1}\) be the discrete Fourier
matrix. Since
\[
  F_q^{-1}P_qF_q=\operatorname{diag}(1,\omega_q,\dots,\omega_q^{q-1}),
\]
conjugating \(I-z\widetilde B_q\) by \(I_{10}\otimes F_q\) gives
\[
  \bigoplus_{j=0}^{q-1}\bigl(I-z(B_0+\omega_q^jB_1)\bigr)
  =\bigoplus_{j=0}^{q-1}\bigl(I-zB(\omega_q^j)\bigr).
\]
Taking determinants yields the stated factorisation.
\end{proof}

\begin{proposition}[Primitive real factors of the cyclic lifts]
  \label{prop:sync-primitive-real-factor}
  Fix \(m\ge 2\). For
  \(a\in(\ZZ/2m\ZZ)^\times/\{\pm1\}\), set
  \[
    r_a:=e^{\pi ia/m},
    \qquad
    s_a:=r_a+r_a^{-1}=2\cos(\pi a/m).
  \]
  Put
  \[
    A_m:=\Bigl\{2\cos\!\tfrac{\pi a}{m}:
      a\in(\ZZ/2m\ZZ)^\times/\{\pm1\}\Bigr\},
    \qquad
    \Phi_m^+(s):=\prod_{\alpha\in A_m}(s-\alpha),
  \]
  and \(d_m:=|A_m|\).
  Then the normalized primitive real factor is the following root
  product, which is the definition of the polynomial used below:
  \[
    R_m(w):=\prod_{\alpha\in A_m}\widehat\Delta(w,\alpha).
  \]
  With the standard Sylvester resultant convention this same polynomial
  is represented by a resultant only in the following order:
  \[
    R_m(w)=
    \Res_s\bigl(\Phi_m^+(s),\widehat\Delta(w,s)\bigr),
  \]
  where \(\Phi_m^+\) is in the first argument. The reversed order is
  not used to define \(R_m\); with the standard Sylvester convention it
  is
  \[
    \Res_s\bigl(\widehat\Delta(w,s),\Phi_m^+(s)\bigr)
    =(-1)^{3d_m}R_m(w).
  \]
  Moreover, for every \(a\) as above,
  \[
    \Delta(z,r_a^2)=\widehat\Delta(zr_a,s_a).
  \]
\end{proposition}

\begin{proof}
The identity
\(\Delta(z,r_a^2)=\widehat\Delta(zr_a,s_a)\) is exactly the defining
change of variables \(u=r^2\), \(w=zr\), \(s=r+r^{-1}\) applied at
\(r=r_a\). If \(a\) is coprime to \(2m\), then \(r_a\) is a primitive
\((2m)\)-th root of unity and \(r_a^2\) is a primitive \(m\)-th root of
unity, so the corresponding factor is one of the primitive Fourier
blocks in \Cref{prop:sync-cyclic-lift-factorisation}. Replacing \(a\) by
\(-a\) conjugates \(r_a\) and leaves \(s_a=r_a+r_a^{-1}\) unchanged, so
conjugate primitive characters collapse to the same real slice
\(\widehat\Delta(w,s_a)\). Thus the displayed product is a polynomial in
the normalized coordinate \(w\): for each primitive character separately,
the corresponding factor in the original cyclic-lift determinant is
obtained by the harmless scaling \(z=w/r_a\).

By definition,
\[
  \Phi_m^+(s)=\prod_{\alpha\in A_m}(s-\alpha),
\]
so \Cref{lem:ordered-cyclotomic-resultant} gives
\[
  \Res_s\bigl(\Phi_m^+(s),\widehat\Delta(w,s)\bigr)
  =
    \prod_{\alpha\in A_m}\widehat\Delta(w,\alpha),
  \qquad d_m=|A_m|.
\]
The root product is therefore represented by a resultant only with
\(\Phi_m^+\) in the first slot. In particular, the notation \(R_m\)
denotes the product itself, not an unordered resultant expression. Since
\(\deg_s\widehat\Delta=3\), the reversed-order identity in
\Cref{lem:ordered-cyclotomic-resultant} gives
\[
    \Res_s\bigl(\widehat\Delta(w,s),\Phi_m^+(s)\bigr)
    =(-1)^{3d_m}
      \Res_s\bigl(\Phi_m^+(s),\widehat\Delta(w,s)\bigr),
\]
which is the displayed reversed-order sign formula.
\end{proof}

\subsection{A general endpoint-selection theorem}

\begin{theorem}[Admissible-cover endpoint selection and asymptotic transfer]
  \label{thm:self-dual-endpoint-selection}
\ifshortcicm
Let \(H(w,s)\in\ZZ[s][w]\) be the completed determinant of a primitive
self-dual labelled kernel with row-sum Perron root \(q\).  The cover
category is the admissible cyclic-cover category of
\Cref{def:admissible-cyclic-cover}: the factors below come only from the
canonical output-cocycle covers \(B_0\otimes I_q+B_1\otimes P_q\), and
not from arbitrary auxiliary one-spike covers.  For
\[
  R_m(w)=\prod_{\alpha\in A_m}H(w,\alpha),
  \qquad
  A_m=\Bigl\{2\cos\!\tfrac{\pi a}{m}:
      a\in(\ZZ/2m\ZZ)^\times/\{\pm1\}\Bigr\},
\]
assume:
\begin{enumerate}[label=\textup{(\roman*)}]
  \item \(H(w,-s)=H(-w,s)\), and \(H(w,2)\) has a simple
        minimal-modulus root \(w_0=1/q\);
  \item the unitary twists have strict spectral gap away from \(s=\pm2\);
  \item on \(2-\eta\le s\le2\) the branch \(w(s)\) through \(w_0\) is
        the unique root of modulus \(<\beta\), where
        \(\beta>q^{-1}\), and is strictly decreasing;
  \item on \(|s|\le2-\eta\), every root has modulus at least
        \(\beta\).
\end{enumerate}
Then for every \(m\) with \(s_m=2\cos(\pi/m)>2-\eta\), the
minimal-modulus roots of \(R_m\) are supplied by
\(H(w,s_m)\) and, when \(m\) is even, by \(H(w,-s_m)\).  Hence
\[
  \rho_m:=\frac{1}{\min\{\abs{w}:R_m(w)=0\}}=\rho(s_m),
  \qquad \rho(s)=1/w(s).
\]
If \(\rho(2-\delta)=\sum_{j=0}^N a_j\delta^j+O(\delta^{N+1})\), then
\[
  \rho_m=\sum_{j=0}^N a_j
    \left(2-2\cos\frac{\pi}{m}\right)^j+O(m^{-2N-2}).
\]
All comparisons with fixed finite-state transfer formulae are over one
fixed coefficient field and one fixed complex embedding; in that
setting every sequence \(u^{\mathsf T}A^m v\) is exponentially bounded,
so factorial or other super-exponential cover budgets cannot be hidden
inside a fixed transfer model.
\else
  Let \(B(u)=B_0+uB_1\) be a finite non-negative output-labelled kernel
  with integer entries. Assume that \(B(1)\) is primitive and that
  \(B(1)\mathbf 1=q\mathbf 1\). Suppose there is a permutation matrix
  \(\Pi\) with
  \[
    \Pi^{-1}B_0\Pi=B_1,\qquad
    \Pi^{-1}B_1\Pi=B_0.
  \]
  Let
  \[
    \Delta(z,u)=\det(I-zB(u)),\qquad
    H(w,s)=\Delta(w/r,r^2),\qquad s=r+r^{-1}.
  \]
  All fixed matrix-transfer statements in this theorem are made over the
  coefficient field \(\QQ\) determined by the displayed integer matrices,
  with the standard embedding \(\QQ\hookrightarrow\CC\).  Thus any
  comparison with a scalar transfer sequence \(u^{\mathsf T}A^m v\) is a
  comparison inside one fixed coefficient field and one fixed complex
  embedding; by \Cref{lem:fixed-field-transfer-exponential-bound}, every
  such fixed transfer sequence is exponentially bounded.
  Then \(H(w,s)\in\ZZ[s][w]\) and \(H(w,-s)=H(-w,s)\). More precisely,
  if \(H(w,s)=\sum_j c_j(s)w^j\), then \(\deg_s c_j\le j\) and
  \(c_j(-s)=(-1)^jc_j(s)\).
  For \(m\ge2\)
  set
  \[
    A_m=\Bigl\{2\cos\!\tfrac{\pi a}{m}:
      a\in(\ZZ/2m\ZZ)^\times/\{\pm1\}\Bigr\},
      \qquad d_m=|A_m|,
  \]
  let
  \[
    \Phi_m^+(s):=\prod_{\alpha\in A_m}(s-\alpha),
  \]
  and define the normalized real cyclotomic factor directly by the root
  product
  \[
    R_m(w):=\prod_{\alpha\in A_m}H(w,\alpha).
  \]
  This definition is made inside the admissible cyclic-cover category of
  \Cref{def:admissible-cyclic-cover}: the characters being grouped are
  the Fourier characters of the canonical covers
  \(B_0\otimes I_q+B_1\otimes P_q\), not characters of an arbitrary
  auxiliary cover.
  This product, not an unordered resultant expression, is the definition
  of \(R_m\). Under the standard Sylvester convention the same
  polynomial is represented by the ordered resultant fixed in
  \Cref{lem:ordered-cyclotomic-resultant},
  \[
    R_m(w)=\Res_s\bigl(\Phi_m^+(s),H(w,s)\bigr).
  \]
  This order is part of the theorem: the cyclotomic polynomial is the
  first argument, so the resultant is the root product itself. The
  notation \(R_m\) never denotes the reversed-order resultant; if the
  two resultant arguments are reversed, the product above is multiplied
  by \((-1)^{(\deg_s H)d_m}\).
  Assume further that
  \[
    \rho(B(e^{it}))<q
    \qquad (t\not\equiv0\pmod{2\pi}),
  \]
  and that \(H(w,2)\) has a simple root \(w_0=1/q\) whose modulus is
  strictly smaller than the moduli of all other roots. Let
  \(w(s)\) be the corresponding analytic branch near \(s=2\), and put
  \(\rho(s)=1/w(s)\). If there are rational constants
  \(\eta>0\) and \(\beta>q^{-1}\) such that
  \[
    \lambda(s):=\min\{\abs{w}:H(w,s)=0\}
  \]
  satisfies
  \[
    \lambda(s)=w(s)<\beta
       \quad (2-\eta\le s\le2),
    \qquad
    \lambda(s)\ge \beta
       \quad (0\le s\le 2-\eta),
  \]
  and \(w(s)\) is strictly decreasing on \([2-\eta,2]\), then for every
  \(m\) with \(s_m=2\cos(\pi/m)>2-\eta\) the root product \(R_m(w)\) has
  its minimal-modulus roots exactly in the factors \(H(w,s_m)\) and,
  when \(m\) is even, \(H(w,-s_m)\). Consequently
  \[
    \rho_m:=\frac{1}{\min\{\abs{w}:R_m(w)=0\}}
    \quad\text{satisfies}\quad
    \rho_m=\rho(s_m).
  \]
  After the displayed hypotheses have been verified for a completed
  determinant \(H\), the proof of this conclusion uses no further
  kernel-specific input.
  Finally suppose that, for some \(N\ge0\),
  \[
    \rho(2-\delta)=\sum_{j=0}^N a_j\delta^j+O(\delta^{N+1})
    \qquad(\delta\to0^+).
  \]
  Then
  \[
    \rho_m
    =\sum_{j=0}^N a_j
      \left(2-2\cos\frac{\pi}{m}\right)^j
      +O(m^{-2N-2}).
  \]
  Equivalently, if
  \[
    2-2\cos\frac{\pi}{m}
      =\sum_{\ell=1}^{N+1}
        \frac{2(-1)^{\ell+1}\pi^{2\ell}}{(2\ell)!}\,m^{-2\ell}
        +O(m^{-2N-4}),
  \]
  then the coefficient of \(m^{-2k}\), \(0\le k\le N\), in \(\rho_m\)
  is
  \[
    \sum_{j=0}^k a_j\,
    [T^k]\left(
      \sum_{\ell=1}^k
        \frac{2(-1)^{\ell+1}\pi^{2\ell}}{(2\ell)!}\,T^\ell
    \right)^j .
  \]
\fi
\end{theorem}

\begin{proof}
\ifshortcicm
For \(\alpha=2\cos\theta\), choose \(r=e^{i\theta}\).  The completed
coordinate change gives
\[
  H(w,\alpha)=\Delta(w/r,r^2),
\]
and \(|r|=1\), so root moduli in \(w\) are the same as root moduli in
the original twisted variable.  The symmetry \(H(w,-s)=H(-w,s)\) gives
\(\lambda(-s)=\lambda(s)\), where
\(\lambda(s)=\min\{|w|:H(w,s)=0\}\).

Among primitive conjugates, \(|2\cos(\pi a/m)|\) is maximal at
\(a=1\), and also at \(a=m-1\) when \(m\) is even.  The monotone
endpoint branch and the central separation therefore make
\(\lambda(s_m)\) strictly smaller than all non-endpoint
\(\lambda(\alpha)\).  Since \(R_m\) is the product of the slices
\(H(w,\alpha)\), its minimal roots come exactly from the endpoint
slice(s).  Substituting
\(\delta_m=2-2\cos(\pi/m)=O(m^{-2})\) into the Taylor expansion of
\(\rho(2-\delta)\) gives the stated asymptotic transfer.
\else
The inclusion \(H\in\ZZ[s][w]\), together with the parity
\(H(w,-s)=H(-w,s)\), is exactly \Cref{lem:self-dual-compression}. The
admissibility clause fixes the cover category before any spectral
selection is made: for each \(q\), the only cover under consideration is
the output-cocycle skew product \(B_0\otimes I_q+B_1\otimes P_q\).
Consequently its character blocks are precisely \(B(\omega_q^j)\), as
in \Cref{prop:sync-cyclic-lift-factorisation}; no other cover data can
create or remove factors in \(R_m\). The
displayed product formula for \(R_m\) is exactly
\Cref{lem:ordered-cyclotomic-resultant} with \(A=A_m\): the monic
polynomial \(\Phi_m^+\) is deliberately placed in the first resultant
slot, so no sign occurs in the root product. The same lemma also gives
the reversed-order sign \((-1)^{(\deg_s H)d_m}\).
For \(\alpha=2\cos\theta\in[-2,2]\)
choose \(r=e^{i\theta}\). Then
\[
  H(w,\alpha)=\Delta(w/r,r^2),
\]
and \(\abs{r}=1\), so the change of variables \(z=w/r\) preserves
moduli. Hence
\[
  \lambda(\alpha)^{-1}=\rho(B(e^{2i\theta})).
  \]
The strict spectral-radius hypothesis gives \(\lambda(s)>1/q\) for
\(-2<s<2\), while \(B(1)\mathbf1=q\mathbf1\) and primitivity give the
unique untwisted Perron root \(w_0=1/q\) at \(s=2\).

The symmetry \(H(w,-s)=H(-w,s)\) gives \(\lambda(-s)=\lambda(s)\).
Thus it is enough to compare \(\abs{\alpha}\). Among the primitive
conjugates
\[
  \alpha=2\cos(\pi a/m),
  \qquad a\in(\ZZ/2m\ZZ)^\times/\{\pm1\},
\]
the maximum of \(\abs{\alpha}\) is \(s_m\), attained only by
\(\alpha=s_m\) and, for even \(m\), by \(\alpha=-s_m\). For all other
conjugates with \(\abs{\alpha}>2-\eta\), strict monotonicity of the
endpoint branch gives \(\lambda(\alpha)=\lambda(\abs{\alpha})
>\lambda(s_m)\). If \(\abs{\alpha}\le2-\eta\), the assumed separation
gives \(\lambda(\alpha)\ge\beta>\lambda(s_m)\). Hence the
minimal-modulus roots of \(R_m(w)=\prod_{\alpha\in A_m}H(w,\alpha)\)
come exactly from the endpoint factor(s), and
\(\rho_m=\rho(s_m)\).

  It remains only to record the coefficient transfer, since this is the
  part used later without reference to the original graph. The only
  fixed-matrix growth fact available in this coefficient-field regime is
  the exponential bound of
  \Cref{lem:fixed-field-transfer-exponential-bound}; in particular, no
  hidden fixed transfer sequence over the stated field and embedding can
  supply factorial growth. Put
  \(\delta_m=2-s_m=2-2\cos(\pi/m)\). The identity just proved gives
\[
  \rho_m=\rho(2-\delta_m).
\]
As \(m\to\infty\), one has \(\delta_m=O(m^{-2})\), and the Taylor
expansion of the cosine gives
\[
  \delta_m
  =\sum_{\ell=1}^{N+1}
    \frac{2(-1)^{\ell+1}\pi^{2\ell}}{(2\ell)!}\,m^{-2\ell}
    +O(m^{-2N-4}).
\]
Substituting this into
\(\rho(2-\delta)=\sum_{j=0}^N a_j\delta^j+O(\delta^{N+1})\) gives
\[
  \rho_m
  =\sum_{j=0}^N a_j\delta_m^j+O(\delta_m^{N+1})
  =\sum_{j=0}^N a_j\delta_m^j+O(m^{-2N-2}).
\]
Expanding the finite polynomial in powers of \(T=m^{-2}\) gives the
displayed coefficient formula. No matrix-theoretic information beyond
the hypotheses of the theorem enters this transfer step.
\fi
\end{proof}

The theorem is the general endpoint-asymptotic result used in the rest
of the paper. For a new self-dual kernel the required inputs are
exactly:
\ifshortcicm
\[
\begin{array}{c|l}
\text{input} & \text{role}\\
\hline
\text{row-sum Perron root and primitivity} & \text{normalises }w_0=1/q\\
\text{self-duality} & \text{gives }H(w,-s)=H(-w,s)\\
\text{unitary spectral gap} & \text{excludes interior endpoint ties}\\
\text{rational interval separation} & \text{makes the endpoint selection effective}\\
\text{endpoint jet} & \text{transfers to the }m^{-2}\text{ expansion}
\end{array}
\]
\else
\begin{enumerate}[label=\textup{(\alph*)}]
  \item a row-sum Perron root \(q\) and primitivity of \(B(1)\);
  \item an involution exchanging the two labelled edge matrices;
  \item the strict unitary spectral gap
  \(\rho(B(e^{it}))<q\) away from \(t=0\);
  \item an endpoint root \(w(2)=1/q\) that is simple and isolated in
  modulus;
  \item an effective interval \(2-\eta\le s\le2\) on which this branch
  is separated from the other roots and is monotone.
\end{enumerate}
Once these verifications are available, no further information about
the original graph is used: the cyclotomic endpoint selection and the
large-\(m\) asymptotic expansion follow from the completed determinant
\(H(w,s)\) alone.
\fi

\subsection{Cyclotomic root products and twisted Perron roots}

\ifshortcicm
For \(m\ge2\), put \(s_m=2\cos(\pi/m)\),
\[
  A_m=\Bigl\{2\cos\!\tfrac{\pi a}{m}:
    a\in(\ZZ/2m\ZZ)^\times/\{\pm1\}\Bigr\},
  \qquad
  \Phi_m^+(s)=\prod_{\alpha\in A_m}(s-\alpha),
\]
and define
\[
  R_m(w)=\prod_{\alpha\in A_m}\widehat\Delta(w,\alpha)
       =\Res_s(\Phi_m^+(s),\widehat\Delta(w,s)).
\]
The resultant order is part of the notation; reversing it multiplies by
\((-1)^{3\deg\Phi_m^+}\).  The twisted Perron root is
\[
  \rho_m=\frac{1}{\min\{\abs{w}:R_m(w)=0\}}.
\]
The case \(m=3\) is exceptional because the leading coefficient
\(s^2-1\) vanishes at \(s=1\), so the degree law starts at \(m=4\).
\else
For \(m\ge 2\), let
\[
  s_m:=2\cos(\pi/m),
\]
let \(\Phi_m^+(s)\) be the minimal polynomial of \(s_m\) over \(\QQ\), and
define the normalized primitive real factor by the conjugate product
\[
  R_m(w):=\prod_{\alpha\in A_m}\widehat\Delta(w,\alpha)\in\ZZ[w],
\]
where \(A_m\) is the conjugate set displayed below. This product is the
definition of \(R_m\). With the standard Sylvester convention the same
polynomial is represented by the ordered resultant
\[
  R_m(w)=\Res_s\bigl(\Phi_m^+(s),\widehat\Delta(w,s)\bigr)
\]
with \(\Phi_m^+\) in the first resultant slot. The order is part of the
notation: since \(\deg_s\widehat\Delta=3\), the same convention gives
\[
  \Res_s\bigl(\widehat\Delta(w,s),\Phi_m^+(s)\bigr)
    =(-1)^{3\deg\Phi_m^+}R_m(w),
\]
so the reversed-order resultant is not used to define \(R_m\).
The notation \(\Phi_m^+\) is non-standard; it denotes the real
cyclotomic polynomial of \(2\cos(\pi/m)\), namely the minimal
polynomial of the real part of a primitive \(2m\)-th root of unity.
Concretely, the conjugates of \(s_m\) under \(\Gal(\QQ(s_m)/\QQ)\) are
\[
  A_m:=\Bigl\{2\cos\!\tfrac{\pi a}{m}:a\in(\ZZ/2m\ZZ)^\times/\{\pm 1\}\Bigr\},
\]
and \(\Phi_m^+(s)=\prod_{\alpha\in A_m}(s-\alpha)\). The degree is
\(\deg\Phi_m^+=|A_m|=\varphi(2m)/2\). By
\Cref{prop:sync-primitive-real-factor}, the polynomial \(R_m(w)\) is
precisely the completed, normalized primitive real factor extracted from
the \(m\)-cyclic lift. It is normalized because each primitive character
\(r_a^2\) uses its own coordinate \(w=zr_a\); before this change of
variables the primitive factors are the ordinary \(z\)-factors
\(\Delta(z,r_a^2)\), and no single unscaled \(z\)-polynomial is being
identified with \(R_m(w)\). Equivalently, the exact identity in the
original and normalized coordinates is
\[
  \Delta(z,r_a^2)=\widehat\Delta(zr_a,s_a),
  \qquad s_a=r_a+r_a^{-1}.
\]
Thus \(R_m\) is a normalized \(w\)-polynomial. It is not literally the
primitive part of \(\det(I-z\widetilde B_m)\) in one common
\(z\)-coordinate, because the rescaling \(w=zr_a\) depends on the
primitive character. With the Sylvester resultant convention, its
ordered-resultant representation is
\[
  R_m(w)=
    \Res_s\bigl(\Phi_m^+(s),\widehat\Delta(w,s)\bigr)
    =
    \prod_{\alpha\in A_m}\widehat\Delta(w,\alpha),
  \qquad d_m:=\deg\Phi_m^+=|A_m|.
\]
Reversing the two resultant arguments would introduce the sign
\((-1)^{3d_m}\); the chosen order avoids that sign in the primitive
factor itself. Thus, with the standard convention,
\[
  \Res_s\bigl(\widehat\Delta(w,s),\Phi_m^+(s)\bigr)
  =(-1)^{3d_m}R_m(w),
\]
and this reversed-order resultant is not used as \(R_m(w)\). For example,
when \(m=2\) one has \(\Phi_2^+(s)=s\) and the two possible orders give
\[
  \Res_s\bigl(s,\widehat\Delta(w,s)\bigr)=\widehat\Delta(w,0),
  \qquad
  \Res_s\bigl(\widehat\Delta(w,s),s\bigr)=-\widehat\Delta(w,0),
\]
which is the sign \((-1)^3\) in the smallest case.

For orientation, the first few real cyclotomic polynomials are
\[
  \Phi_2^+(s)=s,\qquad
  \Phi_3^+(s)=s-1,\qquad
  \Phi_4^+(s)=s^2-2,\qquad
  \Phi_5^+(s)=s^2-s-1.
\]
The cases \(m=2\) and \(m=3\) are exceptional for different reasons:
\(R_2(w)=\widehat\Delta(w,0)\) is handled separately in
\Cref{cor:rho-m2-closed-form}, while for \(m=3\) one has
\(\Phi_3^+(s)=s-1\) and the leading coefficient \(s^2-1\) of
\(\widehat\Delta(w,s)\) vanishes at \(s=1\), so \(\deg R_3=5\) rather
than \(6\). This is why the degree law begins at \(m\ge 4\).

The corresponding twisted Perron root is
\[
  \rho_m:=\frac{1}{\min\{\abs{w}:R_m(w)=0\}}.
\]
\fi

\begin{theorem}[Cyclotomic resultant degree law]
  \label{thm:sync-cyclotomic-degree-law}
  For every \(m\ge 4\),
  \[
    \deg_w R_m=3\varphi(2m).
  \]
  Moreover, if \(m\) is even then \(R_m(w)\) is an even polynomial.
\end{theorem}

\begin{proof}
\ifshortcicm
Here \(\deg\Phi_m^+=|A_m|=\varphi(2m)/2\).  For \(m\ge4\), no
\(\alpha\in A_m\) equals \(\pm1\), so each factor
\(\widehat\Delta(w,\alpha)\) has \(w\)-degree \(6\).  Hence
\(\deg_wR_m=6|A_m|=3\varphi(2m)\).  If \(m\) is even, \(A_m\) is stable
under \(\alpha\mapsto-\alpha\); the sign symmetry
\(\widehat\Delta(w,-s)=\widehat\Delta(-w,s)\) then gives
\(R_m(-w)=R_m(w)\).
\else
Let \(d_m:=\deg\Phi_m^+=\varphi(2m)/2\), and let \(A_m\) be the
conjugate set
\[
  A_m:=\Bigl\{2\cos\!\tfrac{\pi a}{m}:a\in(\ZZ/2m\ZZ)^\times/\{\pm 1\}\Bigr\}.
\]
Then \(\Phi_m^+(s)=\prod_{\alpha\in A_m}(s-\alpha)\), and
\Cref{lem:ordered-cyclotomic-resultant} gives the exact identities
\[
  R_m(w)
  =
    \prod_{\alpha\in A_m}\widehat\Delta(w,\alpha)
  \quad\text{and}\quad
  \Res_s\bigl(\Phi_m^+(s),\widehat\Delta(w,s)\bigr)
  =
    \prod_{\alpha\in A_m}\widehat\Delta(w,\alpha).
\]
The second displayed expression equals the same product. The displayed
resultant order is essential: the
cyclotomic polynomial is in the first slot.
The same lemma gives
\[
  \Res_s\bigl(\widehat\Delta(w,s),\Phi_m^+(s)\bigr)
  =(-1)^{3d_m}R_m(w)
\]
in the reversed order, so the reversed-order resultant is not
identified with \(R_m\).
For \(m\ge 4\), the leading coefficient of \(\widehat\Delta(w,\alpha)\)
as a polynomial in \(w\) is \(\alpha^2-1\). Since \(\alpha=2\cos(\pi a/m)\)
for \(a\in(\ZZ/2m\ZZ)^\times\), we have \(\alpha=\pm 1\) only if
\(\zeta+\zeta^{-1}=\pm 1\) for a primitive \(2m\)-th root of unity
\(\zeta\), which occurs exactly when \(2m=6\), i.e.\ \(m=3\). Hence
for \(m\ge 4\) one has \(\alpha^2-1\neq 0\) for every \(\alpha\in A_m\),
so each factor \(\widehat\Delta(w,\alpha)\) has degree exactly \(6\) in
\(w\). Therefore
\[
  \deg_w R_m=6\,|A_m|=6d_m=3\varphi(2m).
\]

For evenness when \(m\) is even: the set \(A_m\) is stable under
\(\alpha\mapsto-\alpha\), because \(\cos(\pi a/m)\) and
\(-\cos(\pi a/m)=\cos(\pi(m-a)/m)\) belong to the same conjugate set.
The symmetry \(\widehat\Delta(w,-s)=\widehat\Delta(-w,s)\) then gives
\[
\begin{aligned}
  R_m(-w)
  &=
    \prod_{\alpha\in A_m}\widehat\Delta(-w,\alpha) \\
  &=
    \prod_{\alpha\in A_m}\widehat\Delta(w,-\alpha) \\
  &=
    \prod_{\alpha\in A_m}\widehat\Delta(w,\alpha)
    =R_m(w).
\end{aligned}
\]
\fi
\end{proof}

\begin{corollary}[The twist \(m=2\)]
  \label{cor:rho-m2-closed-form}
  One has
\[
    R_2(w)=\widehat\Delta(w,0)
    =1-5w^2+5w^4-w^6
    =(1-w^2)(w^4-4w^2+1),
\]
  and therefore
\[
    \rho_2=\sqrt{2+\sqrt3}.
\]
\end{corollary}

\begin{proof}
Since \(\Phi_2^+(s)=s\), the Sylvester resultant in the chosen order
gives
\[
  R_2(w)=\Res_s(s,\widehat\Delta(w,s))=\widehat\Delta(w,0).
\]
A minimal-modulus zero of \(R_2\) is \(\sqrt{2-\sqrt3}\), whose
reciprocal is the stated value. Together with
\Cref{thm:sync-cyclotomic-degree-law}, this proves
parts~\textup{(i)} and~\textup{(ii)} of
\Cref{thm:intro-kernel-cyclotomic}.
\end{proof}

\subsection{Endpoint jets and large-\texorpdfstring{$m$}{m} asymptotics}

Let \(w(s)\) denote the simple branch of \(\widehat\Delta(w,s)=0\) near
\((w,s)=(1/3,2)\), and set
\[
  \rho(s):=\frac{1}{w(s)}.
\]

\begin{proposition}[Endpoint jet]
  \label{prop:sync-endpoint-jet}
  The branch \(w(s)\) is analytic near \(s=2\), with
  \[
    w(2)=\frac13,\qquad
    w'(2)=-\frac{11}{153},\qquad
    \rho'(2)=\frac{11}{17},\qquad
    \rho''(2)=\frac{1351}{9826}.
  \]
  Moreover, if \(\delta:=2-s\), then
  \begin{align*}
    \rho(2-\delta)
    &=3-\frac{11}{17}\delta+\frac{1351}{19652}\delta^2
      -\frac{17133}{45435424}\delta^3 \\
    &\quad
      -\frac{350037409}{105046700288}\delta^4
      -\frac{47614595705}{242867971065856}\delta^5 \\
    &\quad
      +\frac{141889658001627}{561510749104259072}\delta^6
      +O(\delta^7).
  \end{align*}
\end{proposition}

\begin{proof}
One has
\[
  \widehat\Delta(1/3,2)=0,\qquad
  \partial_w\widehat\Delta(1/3,2)\neq 0.
\]
Hence the implicit function theorem gives a unique analytic branch
\(w(s)\) near \(s=2\). Differentiating the identity
\(\widehat\Delta(w(s),s)=0\) gives
\[
  \partial_w\widehat\Delta(w(s),s)\,w'(s)
  +\partial_s\widehat\Delta(w(s),s)=0.
\]
At \((w,s)=(1/3,2)\) one computes
\[
  \partial_w\widehat\Delta(1/3,2)=-\frac{272}{81},
  \qquad
  \partial_s\widehat\Delta(1/3,2)=-\frac{176}{729},
\]
hence
\[
  w'(2)
  =-\frac{\partial_s\widehat\Delta(1/3,2)}
          {\partial_w\widehat\Delta(1/3,2)}
  =-\frac{11}{153}.
\]
Since \(\rho(s)=w(s)^{-1}\), this immediately gives
\[
  \rho'(2)=-\frac{w'(2)}{w(2)^2}=\frac{11}{17}.
\]
Differentiating a second time yields
\[
  \partial_{ww}\widehat\Delta\,(w')^2
  +2\partial_{ws}\widehat\Delta\,w'
  +\partial_{ss}\widehat\Delta
  +\partial_w\widehat\Delta\,w''=0,
\]
which determines \(w''(2)\), and therefore
\[
  \rho''(2)=\frac{2w'(2)^2}{w(2)^3}-\frac{w''(2)}{w(2)^2}
  =\frac{1351}{9826}.
\]
Equivalently, write
\[
  w(s)=\sum_{j=0}^6 a_j(s-2)^j+O((s-2)^7),
  \qquad a_0=\frac13 .
\]
Put \(t=s-2\), \(w_0=1/3\), and let
\[
  H_n(a_1,\ldots,a_n)
  :=[t^n]\,
    \widehat\Delta\!\left(w_0+\sum_{j=1}^n a_jt^j,\,2+t\right).
\]
Since \(\widehat\Delta(w_0,2)=0\), each equation \(H_n=0\) has the
triangular form
\[
  \partial_w\widehat\Delta(w_0,2)\,a_n
  +H_n(a_1,\ldots,a_{n-1},0)=0,
  \qquad 1\le n\le6.
\]
The coefficient \(\partial_w\widehat\Delta(w_0,2)=-272/81\) is nonzero,
so the coefficients \(a_n\) are determined successively. Solving these
six linear equations gives
\[
\begin{aligned}
  a_1&=-\frac{11}{153},&
  a_2&=\frac{4175}{530604},&
  a_3&=-\frac{336709}{3680269344},\\
  a_4&=\frac{5582818619}{25526348169984},&
  a_5&=-\frac{26151302654713}{177050750907009024},&
  a_6&=\frac{15005008669158719}{1228024008291014590464}.
\end{aligned}
\]
Inverting \(w(s)\) as a formal power series then gives
\[
\begin{aligned}
  \frac{1}{w(2+t)}
  &=3+\frac{11}{17}t+\frac{1351}{19652}t^2
    +\frac{17133}{45435424}t^3
    -\frac{350037409}{105046700288}t^4\\
  &\quad
    +\frac{47614595705}{242867971065856}t^5
    +\frac{141889658001627}{561510749104259072}t^6
    +O(t^7).
\end{aligned}
\]
Setting \(t=-\delta\) yields the displayed jet of
\(\rho(2-\delta)\). The exact recursive identities and the resulting
coefficients displayed in
\Cref{sec:appendix-computation}\,\textup{(f)} are the certificate used
for this endpoint jet.
\end{proof}

\begin{lemma}[Root separation]
  \label{lem:root-separation}
  Put
  \[
    \eta=\frac1{100},\qquad \beta=\frac{101}{300},\qquad m_0=32.
  \]
  For every \(m\ge m_0\), the minimal-modulus roots of
  \(R_m(w)\) are precisely the roots contributed by
  \(\widehat\Delta(w,s_m)\), together with the symmetric factor
  \(\widehat\Delta(w,-s_m)\) when \(-s_m\in A_m\) (equivalently, when
  \(m\) is even). Specifically, letting \(w^*(s_m)\) denote the unique
  root of \(\widehat\Delta(w,s_m)=0\) near \(w=1/3\), one has
  \(\rho_m=\rho(s_m)=1/|w^*(s_m)|\). The remaining cases
  \(4\le m<32\) are verified by exact algebraic interval comparison in
  \Cref{sec:appendix-computation}\,\textup{(e)}.
\end{lemma}

\begin{proof}
\ifshortcicm
Let
\[
  \lambda(s)=\min\{|w|:\widehat\Delta(w,s)=0\}.
\]
The product definition of \(R_m\) gives
\(\rho_m^{-1}=\min_{\alpha\in A_m}\lambda(\alpha)\), and the
sign-symmetry gives \(\lambda(-s)=\lambda(s)\).

At \(s=2\),
\[
  \widehat\Delta(w,2)=(w-1)(w+1)(3w-1)(w^3-w^2+w+1),
\]
so \(w=1/3\) is the unique minimal root.  The implicit-function
calculation of \Cref{prop:sync-endpoint-jet} gives the real branch
\(w^*(s)\) through this point.  The interval certificate printed as
\(\texttt{endpoint\_interval\_certificate: true}\) verifies, on
\(199/100\le s\le2\), that this branch is the unique root of modulus
\(<1/2\), satisfies \(w^*(s)<101/300\), and is strictly decreasing.

For \(|s|\le199/100\), the central interval certificate verifies
\[
  \lambda(s)\ge 101/300.
\]
The strict unitary spectral gap used by the certificate is the standard
triangle-equality argument for the displayed primitive row-stochastic
support: equality in \(\rho(B(e^{it}))\le3\) would force all edge phases
to align along the listed paths, and the rows
\(\st{000},\st{001},\st{010},\st{0-12}\) then force \(e^{it}=1\).

For \(m\ge32\), \(s_m=2\cos(\pi/m)>199/100\).  Among the conjugates in
\(A_m\), the maximum absolute value is attained only by \(s_m\), and
also by \(-s_m\) when \(m\) is even.  The endpoint monotonicity and the
central lower bound therefore make the minimal roots of \(R_m\) exactly
the endpoint slice roots.  The runner separately prints
\(\texttt{finite\_m\_branch\_selection\_4\_to\_31: true}\) for the
remaining finite range.
\else
Define
\[
  \lambda(s):=\min\bigl\{|w|:\widehat\Delta(w,s)=0\bigr\}.
\]
By the symmetry \(\widehat\Delta(w,-s)=\widehat\Delta(-w,s)\), one has
\(\lambda(-s)=\lambda(s)\) for all \(s\). By
\Cref{lem:ordered-cyclotomic-resultant},
\[
  R_m(w)=
    \prod_{\alpha\in A_m}\widehat\Delta(w,\alpha),
  \qquad d_m:=\deg\Phi_m^+,
\]
and \(\rho_m=1/\min\{|w|:R_m(w)=0\}=1/\min_{\alpha\in A_m}\lambda(\alpha)\).
We must show that \(\lambda(s_m)\) is strictly smaller than
\(\lambda(\alpha)\) for all other \(\alpha\in A_m\) when \(m\ge32\).

\medskip
\noindent\textit{Step 1 (Endpoint branch near \(s=2\)).}\quad
At \(s=2\), the polynomial \(\widehat\Delta(w,2)\) factorises as
\[
  \widehat\Delta(w,2)=(w-1)(w+1)(3w-1)(w^3-w^2+w+1).
\]
The minimal-modulus root is \(w=1/3\), and it is simple. The cubic
factor has a unique real root between \(-11/20\) and \(-27/50\) because
it changes sign there. Moreover,
\[
  |w^3-w^2+w|\le |w|^3+|w|^2+|w|<0.54^3+0.54^2+0.54<1
\]
whenever \(|w|\le 27/50\), so \(w^3-w^2+w+1\neq 0\) on the closed disc
of radius \(27/50\). Hence every root other than \(w=1/3\) has modulus
at least \(27/50>1/2\).
By the implicit function theorem, there exist \(\eta>0\) and a unique real-analytic branch
\(w^*:[2-\eta,2]\to\RR\) with \(w^*(2)=1/3\) and
\(\widehat\Delta(w^*(s),s)=0\). Shrinking \(\eta\) if necessary, we may
assume that
\[
  0<w^*(s)<\frac12
  \qquad (2-\eta\le s\le 2),
\]
while every other root of \(\widehat\Delta(\cdot,s)\) has modulus at
least \(1/2\). Hence
\[
  \lambda(s)=w^*(s)
  \qquad (2-\eta\le s\le 2).
\]
This is the same endpoint branch as in \Cref{prop:sync-endpoint-jet},
so \(\rho(s)=1/w^*(s)\) on \([2-\eta,2]\).
Since \(w^{*\prime}(2)=-11/153<0\) by \Cref{prop:sync-endpoint-jet},
continuity of \(w^{*\prime}\) allows us to shrink \(\eta\) further so
that \(w^{*\prime}(s)<0\) on \([2-\eta,2]\). Thus \(w^*(s)>1/3\) for
\(s<2\), and
\[
  \rho(s)=\frac{1}{w^*(s)}
\]
is strictly increasing on \([2-\eta,2]\).
Equivalently \(\lambda(s)=w^*(s)\) is strictly decreasing on this
interval. Together with the symmetry \(\lambda(-s)=\lambda(s)\), this
will compare endpoint-near conjugates through their absolute values.
The exact interval computation in
\Cref{sec:appendix-computation}\,\textup{(e)} verifies that the same
conclusions hold with the explicit interval
\[
  \frac{199}{100}\le s\le2.
\]
More precisely, on this interval the branch \(w^*(s)\) is the unique
root of modulus \(<1/2\), it satisfies
\[
  w^*(s)<\frac{101}{300},
\]
and \(w^{*\prime}(s)<0\). Thus the constants in the statement may be
taken to be \(\eta=1/100\) and \(\beta=101/300\).

\medskip
\noindent\textit{Step 2 (Strict spectral gap away from the untwisted
endpoints).}\quad
Let \(s\in(-2,2)\), choose \(\vartheta\in(0,\pi)\) with
\(s=2\cos\vartheta\), and set \(r=e^{i\vartheta}\), \(u=r^2=e^{2i\vartheta}\).
Then \(\widehat\Delta(w,s)=\Delta(w/r,r^2)\), so
the substitution \(z=w/r\) has \(|z|=|w|\). Hence the minimal
\(|w|\)-root of \(\widehat\Delta(w,s)\) is the minimal \(|z|\)-root of
\(\Delta(z,u)=\det(I-zB(u))\), and \(\lambda(s)^{-1}\) is the spectral
radius of \(B(u)\). This is the comparison between completed
coordinates and the original twisted spectral radius: it uses only that
\(|r|=1\) on the unit circle, so the change of variables preserves
moduli. Since every row
of \(B(1)\) has sum \(3\), one has \(\rho(B(u))\le 3\).

Suppose for contradiction that \(\rho(B(u))=3\). Choose an eigenvalue
\(\mu\) of \(B(u)\) with \(|\mu|=3\) and a corresponding eigenvector
\(\mathbf{x}=(x_1,\dots,x_{10})\), normalised by
\(\max_j |x_j|=1\). Pick \(i\) with \(|x_i|=1\). The \(i\)-th row of
\(B(u)\) contains exactly three non-zero entries, each of modulus \(1\),
so
\[
  3=|\mu x_i|
   =\Bigl|\sum_j b_{ij}(u)x_j\Bigr|
  \le \sum_j |b_{ij}(u)x_j|
  \le 3.
\]
Hence equality holds throughout: the three target coordinates in that
row all have modulus \(1\), and the corresponding summands have a
common phase. Because the digraph is strongly connected, iterating this
argument along directed paths shows that \(|x_j|=1\) for every
\(j\).

Now inspect the specific rows of \(B(u)\). From the row of
\(\st{000}\),
\[
  3=|\mu x_1|=|x_1+x_2+x_3|,
\]
so \(x_1=x_2=x_3\) and \(\mu=3\). Then the row of
\(\st{001}\) gives
\[
  3x_2=x_4+x_5+x_6,
\]
so \(x_4=x_5=x_6=x_1\). The row of \(\st{010}\) reads
\[
  3x_4=x_5+x_6+ux_7.
\]
Since \(x_4=x_5=x_6=x_1\), this forces \(ux_7=x_1\). The row of
\(\st{0-12}\) gives
\[
  3x_7=x_9+x_4+x_5.
\]
All three summands on the right have modulus \(1\), and two of them
are equal to \(x_1\). Equality in the triangle inequality therefore
forces \(x_7=x_9=x_1\). Combining this with \(ux_7=x_1\) gives
\(u=1\), contradicting \(\vartheta\in(0,\pi)\). Therefore
\[
  \lambda(s)>\frac13
  \qquad (-2<s<2).
\]

\medskip
\noindent\textit{Step 3 (Compactness away from the endpoints).}\quad
The function \(\lambda:[-2,2]\to\RR_{>0}\) is continuous. Indeed, since
\(\widehat\Delta(0,s)=1\), its reciprocal polynomial
\[
  t^6\widehat\Delta(1/t,s)
\]
is monic of degree \(6\), so its roots depend continuously on \(s\);
therefore the maximum modulus of its roots, equivalently \(1/\lambda(s)\),
is continuous. On the compact
set
\[
  K:=\left[-\frac{199}{100},\,\frac{199}{100}\right],
\]
Step~2 gives \(\lambda(s)>1/3\) for every \(s\in K\). Hence
\[
  c_0:=\min_{s\in K}\lambda(s)
\]
is well defined and satisfies \(c_0>1/3\).
The interval certificate in
\Cref{sec:appendix-computation}\,\textup{(e)} strengthens this to
\[
  c_0\ge\frac{101}{300}.
\]

\medskip
\noindent\textit{Step 4 (Comparison of endpoint-near conjugates).}\quad
For \(m\ge32\), one has
\[
  s_m=2\cos(\pi/m)>\frac{199}{100}=2-\eta.
\]
Let \(\alpha\in A_m\). If \(|\alpha|>2-\eta\), then
\(|\alpha|\in[2-\eta,2]\). Among the conjugates
\(\alpha=2\cos(\pi a/m)\) with \(1\le a<m\), the quantity \(|\alpha|\)
is maximal at \(a=1\), giving \(s_m=2\cos(\pi/m)\); if \(m\) is even,
the same maximal absolute value also occurs at \(a=m-1\), giving
\(-s_m\). Thus every \(\alpha\in A_m\) with \(\alpha\neq \pm s_m\)
satisfies
\[
  |\alpha|<s_m.
\]
Using \(\lambda(-s)=\lambda(s)\) and the strict decrease of
\(\lambda(s)\) on \([2-\eta,2]\), equivalently the strict increase of
\(\rho(s)=1/\lambda(s)\), we obtain
\[
  \lambda(\alpha)=\lambda(|\alpha|)>\lambda(s_m)
  \qquad
  (\alpha\in A_m,\ \alpha\neq \pm s_m,\ |\alpha|>2-\eta).
\]
If \(|\alpha|\le 2-\eta=199/100\), then \(\alpha\in K\) and Step~3 gives
\[
  \lambda(\alpha)\ge\frac{101}{300}.
\]
The endpoint interval certificate gives
\(\lambda(s_m)<101/300\). Therefore, for every \(m\ge32\),
\[
  \lambda(s_m)\le \lambda(\alpha)
  \qquad
  (\alpha\in A_m),
\]
with strict inequality for \(\alpha\neq \pm s_m\) and equality only for
\(\alpha=-s_m\) when \(-s_m\in A_m\).
Consequently the minimal-modulus roots of \(R_m(w)\) come from the
factor \(\widehat\Delta(w,s_m)\), together with
\(\widehat\Delta(w,-s_m)\) when \(m\) is even, and
\[
  \rho_m=\frac{1}{\lambda(s_m)}=\rho(s_m). \qedhere
\]
\fi
\end{proof}

\begin{corollary}[Endpoint selection for the ten-state kernel]
  \label{cor:sync-endpoint-selection-application}
  The ten-state kernel satisfies the hypotheses of
  \Cref{thm:self-dual-endpoint-selection} with
  \[
    q=3,\qquad
    H(w,s)=\widehat\Delta(w,s),\qquad
    \eta=\frac1{100},\qquad
    \beta=\frac{101}{300}.
  \]
  Consequently, for every \(m\ge32\),
  \[
    \rho_m=\rho\!\left(2\cos\frac{\pi}{m}\right).
  \]
  More generally, any endpoint jet
  \[
    \rho(2-\delta)=\sum_{j=0}^N a_j\delta^j+O(\delta^{N+1})
  \]
  for this branch gives
  \[
    \rho_m
    =\sum_{j=0}^N a_j
      \left(2-2\cos\frac{\pi}{m}\right)^j
      +O(m^{-2N-2}).
  \]
\end{corollary}

\begin{proof}
The row-sum Perron root \(q=3\) and primitivity of \(B(1)\) are
\Cref{prop:sync-kernel-spectrum}. The involution exchanging \(B_0\) and
\(B_1\) is \Cref{prop:sync-self-duality}, and
\Cref{prop:sync-hatdelta-completion} identifies the completed
determinant \(H\) with \(\widehat\Delta\). The endpoint root
\(w(2)=1/3\) is simple and isolated by
\Cref{prop:sync-endpoint-jet} and the endpoint factorisation in the
proof of \Cref{lem:root-separation}. The strict unitary spectral gap and
the effective separation and monotonicity hypotheses with
\(\eta=1/100\) and \(\beta=101/300\) are exactly the assertions
proved in \Cref{lem:root-separation}. Applying
\Cref{thm:self-dual-endpoint-selection} gives both displayed
conclusions.
\end{proof}

\begin{theorem}[Twisted Perron-root asymptotics]
  \label{thm:rho-gap-m12}
  As \(m\to\infty\),
  \begin{align*}
    \rho_m
    &=3-\frac{11\pi^2}{17m^2}
      +\frac{1808\pi^4}{14739m^4}
      -\frac{27872249\pi^6}{2044594080m^6} \\
    &\quad
      -\frac{77645433997\pi^8}{33089710590720m^8}
      +\frac{2976802482983509\pi^{10}}
             {3442653489858508800m^{10}} \\
    &\quad
      +\frac{3832209193654464509\pi^{12}}
             {23878244605658617036800m^{12}}
      +O(m^{-14}).
  \end{align*}
  In particular,
  \[
    3-\rho_m\sim \frac{11\pi^2}{17m^2}.
  \]
\end{theorem}

\begin{proof}
By \Cref{cor:sync-endpoint-selection-application},
\(\rho_m=\rho(s_m)\) for every \(m\ge32\), and this is the range needed
for the asymptotic statement. Thus the branch selection used here is an
application of the general endpoint-selection theorem; the remaining
calculation is the universal coefficient transfer in the explicit
sixth-order jet.
The Taylor series of \(2-s_m=2-2\cos(\pi/m)\) at infinity is
\[
\begin{aligned}
  2-s_m
  &=\frac{\pi^2}{m^2}-\frac{\pi^4}{12m^4}
    +\frac{\pi^6}{360m^6}-\frac{\pi^8}{20160m^8} \\
  &\quad
    +\frac{\pi^{10}}{1814400m^{10}}
    -\frac{\pi^{12}}{239500800m^{12}}+O(m^{-14}).
\end{aligned}
\]
Setting \(\delta=2-s_m\) and substituting into the sixth-order jet of
\Cref{prop:sync-endpoint-jet} yields the displayed asymptotic series,
where the coefficient of \(m^{-2k}\) is determined by the \(\delta^k\)
term of the jet together with the lower-order terms of the Taylor series
of \(\delta\). This proves part~\textup{(iii)} of
\Cref{thm:intro-kernel-cyclotomic}.
\end{proof}

\begin{remark}
\ifshortcicm
The displayed small cases \(m=4,6,8,10\) are checked directly by the
finite certificate key
\(\texttt{finite\_m\_branch\_selection\_4\_to\_31: true}\); the
large-\(m\) asymptotic branch is used only from \(m\ge32\).
\else
The asymptotic expansion is already numerically accurate for moderate
values of \(m\). For the displayed values \(m=4,6,8,10\), the quantity
\(\rho_m\) is verified directly from the definition
\[
  \rho_m^{-1}=\min_{\alpha\in A_m}\min\{\abs{w}:\widehat\Delta(w,\alpha)=0\},
\]
using the finite SageMath checks recorded in
\Cref{sec:appendix-computation}\,\textup{(e)}. Thus the table below does
not use the large-\(m\) part of \Cref{lem:root-separation}. Solving the
specialised equations for these small values at \(100\)-bit precision,
then rounding each displayed decimal to \(12\) digits after the decimal
point and each displayed absolute error to three significant figures,
gives the table below.
The specialised equations used for the positive endpoint slices are
\[
\begin{array}{ll}
m=4:  & \widehat\Delta(w,\sqrt2)=0,\\
m=6:  & \widehat\Delta(w,\sqrt3)=0,\\
m=8:  & \widehat\Delta\bigl(w,\sqrt{2+\sqrt2}\bigr)=0,\\
m=10: & \widehat\Delta\bigl(w,\sqrt{(5+\sqrt5)/2}\bigr)=0.
\end{array}
\]
Comparing the smallest-modulus roots of these equations with the
truncation of
\Cref{thm:rho-gap-m12} gives
\[
  \begin{array}{c|c|c|c}
    m & \rho_m & \text{12th-order approximation} & \text{absolute error} \\
    \hline
    4  & 2.644080451866 & 2.644083856227 & 3.40\times 10^{-6} \\
    6  & 2.831532258235 & 2.831532269934 & 1.17\times 10^{-8} \\
    8  & 2.903081377378 & 2.903081377585 & 2.07\times 10^{-10} \\
    10 & 2.937319429068 & 2.937319429077 & 9.09\times 10^{-12}
  \end{array}
\]
for \(m=4,6,8,10\).
\fi
\end{remark}










